\part{The Anchoring Interface and the Graded Theory}

\secn{18}{18}{Finite experiments and anchoring signatures}

This part defines the anchoring interface and the finite stochastic theory. A coherent family specifies joint channels, including their couplings; repeated looks use distinct acquisition indices. Zero-error support components, equal-law partitions, and Blackwell experiment classes answer different questions. Conditional independence restores the relevant singleton-generated laws for the first two objects, but supplied joint experiments need not be Blackwell joins. Binary-state finite experiments are a lattice; the general larger-state order is not.

\secn{19}{19}{The anchoring interface (H1: signature only)}

The kernel's \hyperlink{stmt:4.2}{Definition~4.2} gives a view as a function \(v : D \to Y_v\). This presupposes---\hspace{0pt}standing assumption (A4)---\hspace{0pt}that every presentation arrives already attributed: it is \emph{of} the unknown candidate, not of something else, and the attribution is correct. Real presentations do not arrive so labeled: an entity mention must be resolved, a sensor trace must be attributed to a source, a data record must be matched to the patient it describes. This section fixes the types in which that attribution step will be modeled. It states no theorems; its purpose is to commit the framework to a signature so that the graded theory of §§20--24, and all later layers, are built against it.

\begin{thmdef}[\hypertarget{stmt:19.1}{Definition~19.1}\ {\normalfont\upshape(raw presentation)}]

A \textbf{raw presentation} is a pair \((y, c)\) with \(y \in Y\) a \textbf{presentation} in a presentation space and \(c \in C\) a \textbf{context} in a context space. The context carries whatever accompanies the presentation without being part of it: provenance metadata, time, channel identity, surrounding discourse.
\end{thmdef}

\begin{thmdef}[\hypertarget{stmt:19.2}{Definition~19.2}\ {\normalfont\upshape(attribution space)}]

An \textbf{attribution space} is a set \(\widehat{T}\) of \textbf{candidate origins}, together with a distinguished element \(\bot \in \widehat{T}\) representing \emph{no admissible origin}. Origins are the things presentations can be \emph{of}: emitting entities, sources, referents.
\end{thmdef}

\begin{thmdef}[\hypertarget{stmt:19.3}{Definition~19.3}\ {\normalfont\upshape(anchoring scenario; origin map)}]

An \textbf{anchoring scenario} for the contrast domain \(D\) is an attribution space \(\widehat{T}\) together with an \textbf{origin map}
\[
o : D \to \widehat{T} \setminus \{\bot\}
\]
assigning to each candidate the origin that presentations emitted under that candidate actually have. Elements of \(\widehat{T} \setminus (o(D) \cup \{\bot\})\) are \textbf{distractors}: admissible-looking origins that no candidate realizes. Two readings of \(o\) delimit the intended range. In the \textbf{tracked-target reading}, \(o\) is constant: \(D\) is the state space of a single observed target, and every presentation genuinely of it shares one origin. In the \textbf{identity reading}, \(o\) is injective: \(D\) is a space of possible identities of the emitter---\hspace{0pt}the entity-resolution regime---\hspace{0pt}and the origin varies with the candidate. Mixed cases are permitted.
\end{thmdef}

\begin{thmdef}[\hypertarget{stmt:19.4}{Definition~19.4}\ {\normalfont\upshape(anchoring correspondence)}]

An \textbf{anchoring correspondence} is a map
\[
\eta : Y \times C \longrightarrow \mathcal{P}(\widehat{T}) \setminus \{\varnothing\},
\]
assigning to each raw presentation its set of admissible attributions. The \textbf{anchoring status} of \((y,c)\) is:

\begin{itemize}
\tightlist
\item
  \textbf{unambiguous} if \(\eta(y,c)\) is a singleton other than \(\{\bot\}\);
\item
  \textbf{ambiguous} if \(|\eta(y,c)| \ge 2\);
\item
  \textbf{failed} if \(\eta(y,c) = \{\bot\}\).
\end{itemize}
\end{thmdef}

\begin{thmdef}[\hypertarget{stmt:19.5}{Definition~19.5}\ {\normalfont\upshape(selection policy; anchoring error)}]

A \textbf{selection policy} for \(\eta\) is a map \(a : Y \times C \to \widehat{T}\) with \(a(y,c) \in \eta(y,c)\): it resolves each raw presentation to a single attribution among the admissible ones. Relative to a scenario \((\widehat{T}, o)\) and a policy \(a\), an anchoring event for a presentation actually emitted under candidate \(Z\) is \textbf{erroneous} (a \textbf{mis-anchoring}) if \(a(y,c) \ne o(Z)\), and \textbf{irreparable} if \(o(Z) \notin \eta(y,c)\)---\hspace{0pt}in the latter case every policy errs, so the defect lies in the correspondence, not in the selection. The distinction matters because the remedies differ: selection error is reducible by a better policy over the same \(\eta\); correspondence error requires changing \(\eta\) itself.
\end{thmdef}

\begin{thmdef}[\hypertarget{stmt:19.6}{Definition~19.6}\ {\normalfont\upshape(anchored view, generalized signature)}]

A \textbf{generalized view} relative to an anchoring scenario \((\widehat{T}, o)\) is a tuple
\[
v = \big( Y_v,\; C_v,\; e_v,\; \eta_v,\; a_v \big)
\]
where \(Y_v\) is a presentation space, \(C_v\) a context space, \(e_v\) an \textbf{emission component} describing how candidates give rise to presentations (in the exact case a map \(e_v : D \to Y_v\); in the graded case of \hyperlink{sec:21}{§21} a channel), \(\eta_v\) an anchoring correspondence on \(Y_v \times C_v\) valued in \(\widehat{T}\), and \(a_v\) a selection policy for \(\eta_v\).
\end{thmdef}

\begin{thmdef}[\hypertarget{stmt:19.7}{Definition~19.7}\ {\normalfont\upshape(kernel recovery; standing assumption A4′)}]

A generalized view is \textbf{soundly and uniquely anchored} if for every \(Z \in D\) and every context \(c \in C_v\) arising with \(e_v(Z)\), the correspondence satisfies
\[
\eta_v\big(e_v(Z), c\big) = \{\,o(Z)\,\}
\]
---\hspace{0pt}unambiguous, unfailed, and correct---\hspace{0pt}whence the selection \(a_v\) is forced and no anchoring event is erroneous. This condition forces the verdict but does not make context uninformative. If context is discarded, collapse to the presentation alone additionally requires \(C\perp Z\mid Y\) (in the deterministic case, context is determined by the presentation). Under that additional condition the generalized view collapses to the kernel view \((Y_v, e_v)\) of \hyperlink{stmt:4.2}{Definition~4.2}; under the tracked-target reading this is exactly the (A4) regime of Part I. Part III adopts, as standing assumption \textbf{(A4′)}, that all views are soundly and uniquely anchored.
\end{thmdef}

\begin{thmrem}[\hypertarget{stmt:19.8}{Remark~19.8}\ {\normalfont\upshape(anchoring is channel-shaped)}]

The interface deliberately types anchoring so that it is not merely a pre-processing stage. The composite \(Z \mapsto o(Z)\) is itself a view on \(D\) in the sense of \hyperlink{stmt:4.2}{Definition~4.2}---\hspace{0pt}the \textbf{identity view}---\hspace{0pt}and the graded versions of the components are channels: emission \(e_v : D \to \Pr(Y_v)\), attribution \(Y_v \times C_v \to \Pr(\widehat{T})\), and their composite a channel \(D \to \Pr(\widehat{T})\) over the same domain as the feature views. Anchoring quality thereby becomes Blackwell-comparable to feature content, and the substantive H1 theory should be expected to \emph{couple} to the graded order of §§20--24 rather than sit orthogonally upstream of it. For this reason the interface fixes types and nothing else: freezing a pipeline (anchor first, then register) would prejudge a coupling that the graded theory is equipped to analyze.
\end{thmrem}

\textbf{What is deliberately not done here.} No dynamics of attribution (how \(\eta\) and \(a\) are computed, from what evidence, at what reliability); no interaction between anchoring and content (how ambiguity, selection error, or irreparable mis-anchoring coarsens, distorts, or corrupts each stratum of content); no graded attribution theory---\hspace{0pt}although every slot of \hyperlink{stmt:19.6}{Definition~19.6} is typed so that grading is a substitution, not a redesign (Remark 19.8). These constitute the substantive H1 theory and are deferred to a later part; the graded machinery about to be built---\hspace{0pt}channels, Blackwell comparison, confusability, sufficiency---\hspace{0pt}is precisely the toolkit that theory will need, which is the reason for the present ordering. Under (A4′), nothing in §§20--24 depends on this section beyond type compatibility.

\secn{20}{20}{Preliminaries for the graded theory}

All material here is standard; sources accompany each notion. The following assumption is in force for the remainder of Part III.

\begin{quote}
\textbf{(A6) Finiteness.} The contrast domain \(D\) and all presentation spaces are finite. Probability distributions are points of the simplex \(\Pr(Y)\); supports are \(\mathrm{supp}\,p = \{y : p(y) > 0\}\).
\end{quote}

\subsecn{1}{20.1}{Channels and garbling}

A \textbf{channel} from \(D\) to \(Y\) is a map \(P : D \to \Pr(Y)\); we write \(P(y \mid Z)\) for the probability of presentation \(y\) given candidate \(Z\). A channel is \textbf{deterministic} if every \(P(\cdot \mid Z)\) is a point mass; deterministic channels are exactly the kernel's views. Channels compose: for \(G : Y \to \Pr(Y')\), the composite \((G \circ P)(y' \mid Z) = \sum_y G(y' \mid y) P(y \mid Z)\).

A channel \(P'\) is a \textbf{garbling} of \(P\), written \(P \succeq_B P'\), if \(P' = G \circ P\) for some channel \(G\). The relation \(\succeq_B\) is a preorder (the \textbf{Blackwell order}); by the Blackwell--Sherman--Stein theorem it coincides, for finite experiments, with uniform superiority in decision problems: \(P \succeq_B P'\) iff for every prior on \(D\) and every finite decision problem the optimal Bayes risk under \(P\) is no worse than under \(P'\) \autocite{blackwell1951,blackwell1953,torgersen1991}. An \textbf{experiment} is a channel from \(D\) considered up to Blackwell equivalence \(\simeq_B\).

\subsecn{2}{20.2}{Sufficiency}

Given a channel \(P : D \to \Pr(Y)\), a map \(T : Y \to Q\) is a \textbf{sufficient statistic} for the family \(\{P(\cdot \mid Z)\}_{Z \in D}\) if the conditional distribution of the presentation given \(T\) is the same for all candidates: for all \(t\) with \(\Pr(T = t \mid Z) > 0\), the distribution \(P(\,\cdot \mid T = t, Z)\) does not depend on \(Z\).

\textbf{\hyperlink{stmt:20.1}{Fact~20.1} (factorization; finite case).} \(T\) is sufficient iff there exist functions \(g, h \ge 0\) with \(P(y \mid Z) = g(T(y), Z)\, h(y)\) for all \(y, Z\). \emph{Proof sketch:} (\(\Leftarrow\)) compute the conditional given \(T = t\) and observe cancellation of \(g(t, Z)\). (\(\Rightarrow\)) take \(h(y) = P(y \mid T = T(y))\) (well defined by sufficiency) and \(g(t, Z) = \Pr(T = t \mid Z)\). \(\square\) (Fisher; \textcite{halmos1949})

A sufficient statistic \(T\) is \textbf{minimal} if it is a function of every sufficient statistic on the realized outcomes. Minimal sufficient statistics exist in the finite case and are constructed in \hyperlink{stmt:22.6}{Theorem~22.6} below \autocite{lehmann1950,bahadur1954}.

\subsecn{3}{20.3}{Information measures}

For jointly distributed finite variables, \(H(\cdot)\), \(H(\cdot \mid \cdot)\), \(I(\cdot\,;\cdot)\), and \(I(\cdot\,;\cdot \mid \cdot)\) denote entropy, conditional entropy, mutual information, and conditional mutual information \autocite{cover2006}. Three standard facts are used repeatedly: \(H(U \mid W) = 0\) iff \(U\) is almost surely a function of \(W\); the \textbf{data processing inequality (DPI)}: if \(Z \to Y \to Y'\) is a Markov chain then \(I(Z; Y') \le I(Z; Y)\), with equality iff \(Z\) and \(Y\) are conditionally independent given \(Y'\) under that joint law (full support of the state prior is needed to conclude sufficiency at every candidate); the \textbf{chain rule}: \(I(Z; Y, Y') = I(Z; Y) + I(Z; Y' \mid Y)\); and \textbf{Fano's inequality}: for any estimator \(\widehat{q}\) of a variable \(q\) with \(M \ge 2\) values, \(\Pr(\widehat q \ne q) \ge \big(H(q \mid Y) - 1\big)/\log_2 M\) \autocite{cover2006}.

\subsecn{4}{20.4}{Zero-error notions}

For a channel \(P\), two candidates are \textbf{confusable} if their output distributions share a support point; the resulting \textbf{confusability graph} underlies Shannon's zero-error theory of channels \autocite{shannon1956}. Zero-error questions of the present framework are governed by this graph (Theorem 22.1), not by any averaged quantity---\hspace{0pt}the first indication that graded content is not one thing.

\subsecn{5}{20.5}{The categorical ambient (remark)}

Everything this part uses---\hspace{0pt}channels and their composition, determinism as a distinguished subclass, marginals, couplings, conditional independence, sufficiency---\hspace{0pt}is native to \textbf{Markov categories}, the synthetic axiomatization of categorical probability \autocite{cho2019,fritz2020}. The development below stays concrete, but three consequences of the categorical reading are worth recording. First, the deterministic/graded relationship is structural: kernel views are the deterministic morphisms of the ambient Markov category, and the recovery theorem (Theorem 24.2) is the statement that the exact theory of Parts I--II is the restriction of the graded theory to that subcategory. Second, categorical treatments of sufficient statistics require their stated structural hypotheses \autocite{fritz2020}. \hyperlink{stmt:22.6}{Theorem~22.6} is proved here for finite experiments; \hyperlink{stmt:A.1}{Proposition~A.1} extends existence of a lossless likelihood statistic to finite states and standard Borel observations. Third, Markov categories provide a language for possible generalization. They do not discharge quotient regularity, uniform identifiability, compactness, or measurable-selection hypotheses; Appendix A states the actual boundary.

\secn{21}{21}{Graded views and the stratification of content}

\subsecn{1}{21.1}{Graded view families}

\begin{thmdef}[\hypertarget{stmt:21.1}{Definition~21.1}\ {\normalfont\upshape(graded view)}]

A \textbf{graded view} on \(D\) is a channel \(v : D \to \Pr(Y_v)\). Under (A4′) its anchoring component is suppressed. Deterministic channels are identified with kernel views.
\end{thmdef}

In the kernel, a set of views determined its compound view outright: the pairing \(\langle v, w \rangle\) was constructed pointwise. Graded views lose this property: the marginal channels of \(v\) and \(w\) do not determine their joint behavior, because the noises may be coupled in different ways. This is a genuinely new degree of freedom with no kernel counterpart, and it must be built into the definition of a family.

\begin{thmdef}[\hypertarget{stmt:21.2}{Definition~21.2}\ {\normalfont\upshape(coherent graded family)}]

A \textbf{coherent graded view family} on \(D\) is a set \(V\) of graded views together with, for each finite \(T \subseteq V\), a \textbf{joint channel} \(P_T : D \to \Pr\big(\prod_{v \in T} Y_v\big)\), such that the assignment is \textbf{projective}: for \(T' \subseteq T\), marginalizing \(P_T\) onto the coordinates of \(T'\) yields \(P_{T'}\); and \(P_{\{v\}} = v\). The family is \textbf{conditionally independent (CI)} if each \(P_T\) is the product of its marginals given \(Z\).
\end{thmdef}

\begin{thmrem}[\hypertarget{stmt:21.3}{Remark~21.3}\ {\normalfont\upshape(coupling indeterminacy)}]

Distinct coherent families can share all their single-view channels and differ in every compound \(P_T\). Content over view sets therefore depends on data the kernel never needed: the coupling. Moreover, a projectively consistent family of \emph{pairwise} joints need not extend to a global joint at all---\hspace{0pt}the classical marginal problem \autocite{vorobev1962}---\hspace{0pt}which is exactly the no-global-section phenomenon analyzed sheaf-theoretically in \autocite{abramsky2011}. Coupling indeterminacy is thus the graded theory's native entry point to hook H4 and is flagged here for that purpose. It is convenient to name the object: the \textbf{coupling fiber} over a marginal system \(\{v\}_{v \in V}\) is the class of coherent families extending it. Vorob'ev's theorem characterizes, in the exact case, the overlap structures over which the fiber is never empty; content over view sets is a function on the fiber, constant across it exactly when the coupling mechanism of \hyperlink{stmt:22.4}{Proposition~22.4}(2) is absent. (The flag is cashed in Part V, \hyperlink{sec:39}{§39}: Vorob'ev's theorem is the unconditional-descent criterion of the coupling presheaf, and the fiber named here carries both of H4's failure modes---\hspace{0pt}contextual evidence as its emptiness, underdetermination as its multiplicity---\hspace{0pt}with coupling superadditivity located as non-constancy of \(\sigma^{=}\) on fibers; \hyperlink{stmt:39.2}{Theorem~39.2}, \hyperlink{stmt:39.3}{Proposition~39.3}.)
\end{thmrem}

\subsecn{2}{21.2}{Graded content: three strata}

The kernel had one content object, \(\sigma(S) \in \mathrm{Part}(D)\). Grading splits it into three, ordered by how much of the channel structure they retain.

\begin{thmdef}[\hypertarget{stmt:21.4}{Definition~21.4}\ {\normalfont\upshape(graded content, full stratum)}]

For finite \(S \subseteq V\) of a coherent family, the \textbf{graded content} is the experiment
\[
\widehat{\sigma}(S) \;=\; [\,P_S\,]_{\simeq_B},
\]
the Blackwell class of the joint channel.
\end{thmdef}

\begin{thmdef}[\hypertarget{stmt:21.5}{Definition~21.5}\ {\normalfont\upshape(statistical stratum)}]

The \textbf{graded registered representation} is the likelihood map
\[
\widehat{R}_S : D \to \Pr\Big(\textstyle\prod_{v \in S} Y_v\Big), \qquad \widehat{R}_S(Z) = P_S(\cdot \mid Z),
\]
and the \textbf{statistical content} is its kernel partition
\[
\sigma^{=}(S) \;=\; \ker\big(\widehat{R}_S\big) \in \mathrm{Part}(D):
\]
candidates are identified iff their output distributions are identical.
\end{thmdef}

\begin{thmdef}[\hypertarget{stmt:21.6}{Definition~21.6}\ {\normalfont\upshape(zero-error stratum)}]

The \textbf{confusability relation} on \(D\) is
\[
Z \approx_S Z' \;:\iff\; \mathrm{supp}\, P_S(\cdot \mid Z) \,\cap\, \mathrm{supp}\, P_S(\cdot \mid Z') \;\ne\; \varnothing,
\]
a reflexive and symmetric relation that is \textbf{not} transitive in general. The \textbf{zero-error content} \(\sigma^{0}(S) \in \mathrm{Part}(D)\) is the partition of \(D\) into connected components of \(\approx_S\).
\end{thmdef}

\hyperlink{stmt:21.5}{Definition~21.5} is the graded \hyperlink{stmt:5.1}{Definition~5.1}--5.2: the registered representation of a candidate is no longer the set of unexcluded alternatives but the likelihood it induces, and indistinguishability is equality of likelihoods. \hyperlink{stmt:21.6}{Definition~21.6} captures a strictly more demanding operational notion---\hspace{0pt}single-observation certainty---\hspace{0pt}and its non-transitivity is the first structural novelty of grading: exact indistinguishability was an equivalence; confusability is only a tolerance, and content extraction must pass to components.

\begin{thmplain}[\hypertarget{stmt:21.7}{Lemma~21.7}\ {\normalfont\upshape(the strata are ordered)}]

For every finite \(S\):
\[
\sigma^{0}(S) \;\le\; \sigma^{=}(S) \qquad \text{in } \mathrm{Part}(D),
\]
and both are monotone in \(S\): \(T \subseteq S\) implies \(\sigma^{0}(T) \le \sigma^{0}(S)\) and \(\sigma^{=}(T) \le \sigma^{=}(S)\). In the deterministic case \(\sigma^{0}(S) = \sigma^{=}(S) = \sigma(S)\).
\end{thmplain}

\begin{proof}

If \(\widehat R_S(Z) = \widehat R_S(Z')\) the two supports are equal, hence intersect (they are nonempty), so \(Z \approx_S Z'\) and the two candidates share a \(\sigma^0\)-component; thus every \(\sigma^=\)-block lies in a \(\sigma^0\)-block, i.e.~\(\sigma^0(S) \le \sigma^=(S)\). Monotonicity of \(\sigma^=\): if the joint likelihoods over \(S\) are equal, their marginals over \(T\) are equal. Monotonicity of \(\sigma^0\): if joint supports over \(S\) intersect at a point, its coordinates witness intersection of the marginal supports over \(T\); so \(\approx_S\; \subseteq\; \approx_T\) as edge sets, hence components over \(S\) refine components over \(T\). Deterministic case: point-mass likelihoods are equal iff the presentations are equal, and supports intersect iff the presentations are equal, so both relations coincide with \(\sim_S\) and both partitions with \(\sigma(S)\), the confusability relation being already transitive.
\end{proof}

The interpretive content of the strata is fixed by the answerability theorems of the next section: \(\sigma^0\) governs what can be answered \emph{with certainty from one look}, \(\sigma^=\) governs what can be answered \emph{with vanishing error given unlimited repetition}, and \(\widehat\sigma\) retains everything, including all intermediate error/prior trade-offs.

\secn{22}{22}{Fate of the kernel theorems under grading}

\subsecn{1}{22.1}{Answerability grades into a family of notions}

\begin{thmdef}[\hypertarget{stmt:22.0}{Definition~22.0}]

Fix a question \(q : D \to A_q\) and finite \(S\). A (deterministic) \textbf{decision rule} is a map \(\alpha : \prod_{v\in S} Y_v \to A_q\). The question is:

\begin{itemize}
\tightlist
\item
  \textbf{zero-error answerable from \(S\)} if some rule satisfies \(\Pr\big(\alpha(Y_S) = q(Z) \,\big|\, Z\big) = 1\) for every \(Z \in D\);
\item
  \textbf{\(\varepsilon\)-answerable from \(S\) under prior \(\mu\)} if some rule has Bayes error \(e_\mu(q \mid S) := \Pr_{\mu}\big(\alpha(Y_S) \ne q(Z)\big) \le \varepsilon\);
\item
  \textbf{asymptotically answerable from \(S\)} if, with \(n\) conditionally i.i.d. copies of \(Y_S\) given \(Z\), there are rules \(\alpha_n\) whose worst-case error \(\max_{Z} \Pr(\alpha_n \ne q(Z) \mid Z)\) tends to \(0\).
\end{itemize}
\end{thmdef}

\begin{thmplain}[\hypertarget{stmt:22.1}{Theorem~22.1}\ {\normalfont\upshape(zero-error answerability; fate of Proposition 5.6 and
Theorem 6.4)}]

\(q\) is zero-error answerable from \(S\) iff \(\ker(q) \le \sigma^{0}(S)\). Consequently the zero-error answerable questions form the principal ideal \({\downarrow}\,\sigma^{0}(S)\), and \(\sigma^0(S)\), as a question, is the maximally informative zero-error answerable one.
\end{thmplain}

\begin{proof}

(\(\Rightarrow\)) Let \(\alpha\) be a perfect rule. For every \(Z\) and every \(y \in \mathrm{supp}\,P_S(\cdot \mid Z)\), correctness with probability one over a finite support forces \(\alpha(y) = q(Z)\). If \(Z \approx_S Z'\), a common support point \(y\) gives \(q(Z) = \alpha(y) = q(Z')\). So \(q\) is constant on \(\approx_S\)-edges, hence on connected components, i.e.~\(\ker(q) \le \sigma^0(S)\). (\(\Leftarrow\)) Suppose \(q\) is constant on components. For each realized \(y\) (in some support) set \(\alpha(y) = q(Z_y)\) for any \(Z_y\) whose support contains \(y\); the value is well defined because any two such candidates are \(\approx_S\)-related. For arbitrary \(Z\) and any \(y \in \mathrm{supp}\,P_S(\cdot\mid Z)\) we get \(\alpha(y) = q(Z)\), so the rule is perfect. The ideal statement follows as in \hyperlink{stmt:6.4}{Theorem~6.4}.
\end{proof}

\begin{thmplain}[\hypertarget{stmt:22.2}{Theorem~22.2}\ {\normalfont\upshape(asymptotic answerability)}]

\(q\) is asymptotically answerable from \(S\) iff \(\ker(q) \le \sigma^{=}(S)\).
\end{thmplain}

\begin{proof}

(\(\Leftarrow\)) There are finitely many distinct likelihoods \(\{\widehat R_S(Z)\}\); let \(\delta > 0\) be the minimum total-variation distance between distinct ones. Let \(\widehat p_n\) be the empirical distribution of the \(n\) copies; define \(\alpha_n\) to output \(q(Z^*)\) for any \(Z^*\) minimizing \(\|\widehat p_n - \widehat R_S(Z^*)\|_{TV}\). Given \(Z\), the law of large numbers gives \(\|\widehat p_n - \widehat R_S(Z)\| \to 0\) a.s., so eventually the minimizer's likelihood equals \(\widehat R_S(Z)\), i.e.~\(Z^* \sim_{\sigma^=} Z\); since \(q\) is constant on \(\sigma^=\)-blocks, \(\alpha_n\) is eventually correct, and finiteness of \(D\) upgrades this to uniform error \(\to 0\) (standard concentration, e.g.~via Hoeffding on each cell). (\(\Rightarrow\)) If \(\widehat R_S(Z) = \widehat R_S(Z')\) with \(q(Z) \ne q(Z')\), then for any rule \(\alpha_n\) the two error probabilities are computed under the \emph{same} output distribution, and since \(\alpha_n\) cannot equal both values at once, \(\Pr(\alpha_n \ne q(Z) \mid Z) + \Pr(\alpha_n \ne q(Z') \mid Z') \ge 1\); the worst-case error is \(\ge 1/2\) for every \(n\).
\end{proof}

\hyperlink{stmt:22.1}{Theorems~22.1}--22.2 give the strata their meaning and show the kernel's single answerability criterion (Proposition 5.6) splits into inequivalent graded criteria with the \emph{same lattice form}: factorization below a content partition. By contrast, \(\varepsilon\)-answerability at fixed \(\varepsilon \in (0, \tfrac12)\) is a down-set condition on \(\ker(q)\) (coarsening a question cannot increase its Bayes error: post-compose the optimal rule) but \textbf{not} a principal ideal in general---\hspace{0pt}the same non-principality already met in the budget-type admissibility of \hyperlink{sec:16.2}{§16.2}, now arising intrinsically. Both instances, and a third at \hyperlink{stmt:23.1}{Definition~23.1}, are consolidated by the ideal completion of \hyperlink{sec:24}{§24} (Remark 24.6).

\subsecn{2}{22.2}{Fate of T1: the adjunction breaks twice}

\begin{thmplain}[\hypertarget{stmt:22.3}{Proposition~22.3}\ {\normalfont\upshape(zero-error content is superadditive)}]

For finite \(S, S'\) in a coherent family,
\[
\sigma^{0}(S \cup S') \;\ge\; \sigma^{0}(S) \vee \sigma^{0}(S'),
\]
and the inequality can be strict even for CI families. Hence \(\sigma^0\) does not preserve joins and admits no upper adjoint; the adjunction of \hyperlink{stmt:6.2}{Theorem~6.2} fails for zero-error content.
\end{thmplain}

\begin{proof}

Inequality: if joint supports over \(S \cup S'\) intersect at a point, its coordinate projections witness intersection over \(S\) and over \(S'\); so \(\approx_{S\cup S'} \subseteq \approx_S \cap \approx_{S'}\) as edge sets, and every component of the smaller edge set is contained in a component of each larger one, hence in a block of \(\sigma^0(S) \vee \sigma^0(S')\). Strictness: let \(D = \{1,2,3\}\) and take CI views \(v, w\) with supports
\[
\mathrm{supp}\,v(\cdot|1)=\{a\},\quad \mathrm{supp}\,v(\cdot|2)=\{a,b\},\quad \mathrm{supp}\,v(\cdot|3)=\{b\};
\]
\[
\mathrm{supp}\,w(\cdot|1)=\{c,d\},\quad \mathrm{supp}\,w(\cdot|2)=\{c\},\quad \mathrm{supp}\,w(\cdot|3)=\{d\}
\]
(any strictly positive probabilities on these supports). Then \(\approx_{\{v\}}\) has edges \(\{12, 23\}\) and \(\approx_{\{w\}}\) has edges \(\{12, 13\}\); each graph is connected, so \(\sigma^0(\{v\}) = \sigma^0(\{w\}) = \bot\) and the join is \(\bot\). Under CI the joint support is the product of supports, so \(\approx_{\{v,w\}}\) has edge set \(\{12,23\} \cap \{12,13\} = \{12\}\), whence \(\sigma^0(\{v,w\}) = \{\,12 \mid 3\,\} > \bot\).
\end{proof}

Two views, each singly incapable of certifying anything with certainty, jointly certify a real distinction. This is a \emph{noise-induced} analogue of Part II's registration anomaly: superadditivity of content arising not from a restricted decoder but from the support geometry of the channels themselves. It is the second independent mechanism producing ``the whole reads more than its parts.''

The mechanism can, however, be located with complete precision, and doing so shows that the zero-error stratum has a lossless invariant on which compositionality survives.

\begin{thmplain}[\hypertarget{stmt:22.3′}{Proposition~22.3′}\ {\normalfont\upshape(the confusability graph is the lossless zero-error
object)}]

Let \(\mathrm{Tol}(D)\) be the set of reflexive symmetric relations (\textbf{tolerances}) on \(D\)---\hspace{0pt}a complete lattice under inclusion of edge sets---\hspace{0pt}and order it by \emph{reverse} inclusion, written \(\sqsubseteq\) (fewer confusions \(=\) more information), under which it is again a complete lattice with joins given by intersection of edge sets. Define the \textbf{graph-valued zero-error content} of finite \(S\) as
\[
\sigma^{G}(S) \;:=\; \approx_S \;\in\; \mathrm{Tol}(D).
\]
Then, on CI families,
\[\begin{gathered}\sigma^{G}(S) \;=\; \bigsqcup_{v \in S} \sigma^{G}(\{v\}) \\ \big(\text{equivalently } \approx_S \,=\, \textstyle\bigcap_{v \in S} \approx_{\{v\}} \text{ as edge sets}\big),\end{gathered}\]
so \(\sigma^{G}\) preserves unions as joins, and, for finite \(V\) (or the arbitrary-intersection extension to all subsets), the full adjunction of \hyperlink{stmt:6.2}{Theorem~6.2} holds for graph-valued zero-error content over CI families, with generators \(\sigma^{G}(\{v\})\) and upper adjoint \(\tau^{G}(\theta) = \{v \in V : \sigma^{G}(\{v\}) \sqsubseteq \theta\}\) (in edge sets: every confusion \(\theta\) permits is one the view permits---\hspace{0pt}the views individually no more informative than \(\theta\)).
\end{thmplain}

\begin{proof}

Under CI, \(\mathrm{supp}\, P_S(\cdot \mid Z) = \prod_{v \in S} \mathrm{supp}\, v(\cdot \mid Z)\), and products of nonempty sets intersect iff every pair of corresponding factors intersects; so \(Z \approx_S Z'\) iff \(Z \approx_{\{v\}} Z'\) for every \(v \in S\), which is the displayed join law. The adjunction then follows exactly as in \hyperlink{stmt:15.1}{Theorem~15.1}: any content map given as a join of singleton generators into a complete lattice is a left adjoint.
\end{proof}

\textbf{Mathematical scope.} Here ``lossless'' refers to retaining the confusability relation, not the full experiment. With the reverse-edge order, component formation is a \textbf{right} adjoint to inclusion of equivalence relations: \(\pi\sqsubseteq R\iff\pi\le\operatorname{comp}(R)\). Its failure to preserve joins is therefore unsurprising. ``Reflection'' remains the manuscript's descriptive mechanism name, not the orientation of this adjunction.

\hyperlink{stmt:22.3′}{Proposition~22.3′} diagnoses \hyperlink{stmt:22.3}{Proposition~22.3} exactly. The partition \(\sigma^{0}(S)\) is the image of \(\sigma^{G}(S)\) under the reflection of tolerances onto equivalence relations (transitive closure; connected components), and the superadditivity of \(\sigma^{0}\) on CI families is precisely the failure of that reflection to commute with edge-intersection---\hspace{0pt}an order-theoretic fact about a closure operator, not a fact about channels. At the graph level, zero-error content over CI families is strictly compositional; the mechanism of this second anomaly is thus better named \textbf{reflection}, with support geometry as its carrier, and off the CI class the joint support is no longer a product, so the law fails through coupling---\hspace{0pt}consistently with the mechanism ladder below. Taking the graph as the primary invariant also aligns the stratum with the zero-error literature, where the confusability graph, not its component partition, carries the theory \autocite{shannon1956,lovasz1979}.

The statistical stratum sits strictly between the zero-error stratum and the full one, and its union behavior completes the picture:

\begin{thmplain}[\hypertarget{stmt:22.4}{Proposition~22.4}\ {\normalfont\upshape(statistical content: superadditivity exactly through
coupling)}]

For finite \(S, S'\) in a coherent family:

\begin{enumerate}
\def\labelenumi{\arabic{enumi}.}
\tightlist
\item
  \(\sigma^{=}(S \cup S') \;\ge\; \sigma^{=}(S) \vee \sigma^{=}(S')\);
\item
  the inequality can be strict, and the mechanism is coupling: candidates can have identical marginal likelihoods over each of \(S\) and \(S'\) yet distinct joint likelihoods over \(S \cup S'\);
\item
  on CI families, \(\sigma^{=}(S) = \bigvee_{v \in S} \sigma^{=}(\{v\})\) for every finite \(S\); hence \(\sigma^{=}\) preserves unions as joins, and the full adjunction of \hyperlink{stmt:6.2}{Theorem~6.2} holds for statistical content, with generators \(\sigma^{=}(\{v\})\).
\end{enumerate}
\end{thmplain}

\begin{proof}

1: Equal joint likelihoods have equal marginals (projectivity), so every \(\sigma^{=}(S \cup S')\)-block lies inside a block of \(\sigma^{=}(S)\) and of \(\sigma^{=}(S')\). 2: Let \(D = \{Z, Z'\}\) and let \(v, w\) be binary views whose four marginal likelihoods are all uniform on \(\{0,1\}\), so \(\sigma^{=}(\{v\}) = \sigma^{=}(\{w\}) = \bot\); couple the noises so that under \(Z\) the two outputs are equal with probability one, while under \(Z'\) they are independent. Both joints have uniform marginals, so the family is projective; the joint likelihoods differ (one is supported on the diagonal), so \(\sigma^{=}(\{v,w\}) = \top > \bot = \sigma^{=}(\{v\}) \vee \sigma^{=}(\{w\})\). 3: Under CI, \(P_S(\cdot \mid Z) = \bigotimes_{v \in S} v(\cdot \mid Z)\), and a product distribution determines and is determined by its factors; hence \(\widehat{R}_S(Z) = \widehat{R}_S(Z')\) iff \(v(\cdot \mid Z) = v(\cdot \mid Z')\) for every \(v \in S\), i.e.~\(\ker \widehat{R}_S = \bigvee_{v \in S} \ker \widehat{R}_{\{v\}}\) as in (3.1). The adjunction follows as in \hyperlink{stmt:6.2}{Theorem~6.2}, the content map being again a join over singleton generators.
\end{proof}

At the full stratum the failure is even more basic:

\begin{thmplain}[\hypertarget{stmt:22.5}{Proposition~22.5}\ {\normalfont\upshape(joint experiments and Blackwell joins)}]

The graded content map is monotone: marginalizing \(P_S\) to \(T\subseteq S\) is a garbling. Thus \(P_{S\cup S'}\) is an upper bound of \(P_S,P_{S'}\). It need not be their least upper bound: marginal channels do not determine their coupling, and the chosen coupling can add information. Finite binary-state experiments do have Blackwell joins; for larger state spaces joins can fail \autocite{bertschinger2014}. Deterministic experiments have unique joint coupling and their join is the common-refinement partition. Appendix D.5 characterizes least members of full coupling fibers.
\end{thmplain}

\begin{proof}

Projections prove the order assertions. The equal-versus-independent binary coupling example of \hyperlink{stmt:22.4}{Proposition~22.4} gives different joint information for fixed null marginals. Deterministic marginals force a paired point mass, and \hyperlink{stmt:24.1}{Lemma~24.1} identifies its least-upper-bound property. The general existence/nonexistence distinction is the cited Blackwell lattice theorem.
\end{proof}

The fate of T1 is thus a trichotomy across the strata, ordered by how much structure fails and why. Zero-error content is superadditive even for CI families---\hspace{0pt}Proposition 22.3's example is CI---\hspace{0pt}with support geometry as the carrier and reflection as the mechanism: at the graph level the CI composition law is exact, and it is the passage to components that destroys it (Proposition 22.3′). Statistical content is superadditive exactly through coupling, and on CI families recovers the full adjunction: at this stratum the kernel's compositionality is not lost but \emph{conditional}. Full graded content loses the very lattice in which a join law could be stated, and its unions acquire an input---\hspace{0pt}the coupling---\hspace{0pt}that the parts do not determine. Under admissibility (Part II) the adjunction failed \emph{contingently}, at non-separable structures; under grading it degrades \emph{structurally}, and the responsible mechanism migrates upward: from reflection (of the confusability graph onto partitions), to coupling, to the order itself. The kernel is the locus where all three pathologies are invisible.

\subsecn{3}{22.3}{Fate of T4: minimal sufficiency}

The kernel's \hyperlink{stmt:6.7}{Theorem~6.7} characterized \(D/\sigma(S)\) as the terminal lossless consolidation. Grading relocates the consolidation from the state side to the evidence side---\hspace{0pt}one compresses the outcome space of the experiment rather than the domain---\hspace{0pt}and the terminal object becomes the minimal sufficient statistic.

\begin{thmplain}[\hypertarget{stmt:22.6}{Theorem~22.6}\ {\normalfont\upshape(graded T4)}]

Let \(P_S\) be the joint channel of a finite \(S\) and let \(Y^+ \subseteq \prod_{v\in S} Y_v\) be its realized outcomes (\(y \in Y^+\) iff \(P_S(y \mid Z) > 0\) for some \(Z\)). Define on \(Y^+\):
\[
y \equiv y' \;:\iff\; \exists\, c > 0\ \ \forall Z \in D:\ \ P_S(y \mid Z) = c\, P_S(y' \mid Z)
\]
(proportional likelihood profiles), and let \(m : Y^+ \to Y^+/{\equiv}\) be the quotient map. Then:

\begin{enumerate}
\def\labelenumi{\arabic{enumi}.}
\tightlist
\item
  \(m\) is a sufficient statistic for \(\{P_S(\cdot \mid Z)\}_Z\);
\item
  \(m\) is minimal: for every sufficient statistic \(T\) there is a map \(f\) with \(m = f \circ T\) on \(Y^+\);
\item
  sufficiency is exactly losslessness: a statistic \(T\) is sufficient iff \(I(Z; T(Y_S)) = I(Z; Y_S)\) for every prior on \(D\) (equality in the DPI);
\item
  in the deterministic case, \(Y^+ = \mathrm{Im}\,\langle v \rangle_{v \in S}\), the classes of \(\equiv\) are singletons, and \(m\) is (up to iso) the identity on \(\mathrm{Im}\,\langle v\rangle_{v\in S} \cong D/\sigma(S)\): \hyperlink{stmt:6.7}{Theorem~6.7}'s terminal consolidation is recovered exactly.
\end{enumerate}
\end{thmplain}

\begin{proof}

1: Fix a representative \(y_0\) in each class; for \(y \equiv y_0\) write \(P_S(y \mid Z) = c_y P_S(y_0 \mid Z)\). Then \(P_S(y \mid Z) = g(m(y), Z)\, h(y)\) with \(g([y_0], Z) = P_S(y_0 \mid Z)\) and \(h(y) = c_y\); sufficiency by \hyperlink{stmt:20.1}{Fact~20.1}. 2: Let \(T\) be sufficient with factorization \(P_S(y \mid Z) = g(T(y), Z) h(y)\); for realized \(y\), \(h(y) > 0\). If \(T(y) = T(y') = t\) with \(y, y' \in Y^+\), the factorization gives \(P_S(y \mid Z) = g(t, Z)\, h(y)\) and \(P_S(y' \mid Z) = g(t, Z)\, h(y')\) for every \(Z\) simultaneously, whence \(P_S(y \mid Z) = \tfrac{h(y)}{h(y')}\, P_S(y' \mid Z)\) with \(\tfrac{h(y)}{h(y')} > 0\)---\hspace{0pt}no ratio of possibly-zero likelihoods is ever formed---\hspace{0pt}so \(y \equiv y'\) and \(m(y) = m(y')\); hence \(m\) factors through \(T\) on \(Y^+\). 3: Equality in the DPI for all priors iff the Markov chain \(Z \to T(Y_S) \to Y_S\) holds in the conditional sense defining sufficiency \autocite{cover2006}. 4: With point-mass likelihoods, \(P_S(y \mid Z) = \mathbb{1}\{\langle v\rangle(Z) = y\}\); two distinct realized outcomes have disjoint, nonempty preimages, so their profiles are never proportional; \(\equiv\) is equality, and \(\mathrm{Im}\,\langle v \rangle \cong D/\sigma(S)\) is \hyperlink{stmt:6.7}{Theorem~6.7}(3).
\end{proof}

\hyperlink{stmt:22.6}{Theorem~22.6} also grades \hyperlink{stmt:4.7}{Lemma~4.7}. In the kernel, sound registrations were bracketed between constants and the canonical (most informative) one; in the graded theory, \emph{lossless} registrations of an experiment are bracketed between the identity and the minimal sufficient statistic (the coarsest lossless one, by items 2--3), and any coarsening past \(m\) begins to lose content, with the DPI quantifying the loss. Registration theory thus acquires a floor with a classical name.

\subsecn{4}{22.4}{The Shannon bridge}

Fix a coherent family, finite \(S\), a question \(q\), and a \textbf{full-support prior} \(\mu\) on \(D\); let \(Z \sim \mu\) and \(Y_S \sim P_S(\cdot \mid Z)\).

\begin{thmplain}[\hypertarget{stmt:22.7}{Theorem~22.7}\ {\normalfont\upshape(entropic criteria)}]

\par\nopagebreak

\begin{enumerate}
\def\labelenumi{\arabic{enumi}.}
\tightlist
\item
  \(H\big(q(Z) \mid Y_S\big) = 0\) for one (equivalently, every) full-support prior iff \(\ker(q) \le \sigma^{0}(S)\)---\hspace{0pt}the vanishing of conditional entropy is exactly zero-error answerability.
\item
  (Fano, quantitative unanswerability.) For any rule, with \(M = |A_q| \ge 2\):
  \[
  \Pr\big(\alpha(Y_S) \ne q(Z)\big) \;\ge\; \frac{H(q(Z) \mid Y_S) - 1}{\log_2 M}.
  \]
\item
  (Chain rule; graded T3.) \(I\big(Z; Y_{S \cup \{w\}}\big) = I(Z; Y_S) + I(Z; Y_w \mid Y_S) \ge I(Z;Y_S)\), and \(e_\mu(q \mid S \cup \{w\}) \le e_\mu(q \mid S)\): a further view never hurts, and its marginal value is the conditional mutual information, null iff \(Y_w \perp Z \mid Y_S\).
\item
  (DPI; graded \hyperlink{stmt:14.2}{Lemma~14.2}.) For any registration channel \(\kappa\) applied to \(Y_S\): \(I(Z; \kappa(Y_S)) \le I(Z; Y_S)\)---\hspace{0pt}registration cannot create information, with equality exactly at sufficiency.
\end{enumerate}
\end{thmplain}

\begin{proof}

1: \(H(q(Z) \mid Y_S) = 0\) iff for every realized \(y\), \(q\) is constant on \(\{Z : \mu(Z) P_S(y \mid Z) > 0\}\); with \(\mu\) full-support these are the sets \(\{Z : y \in \mathrm{supp}\, P_S(\cdot\mid Z)\}\). Constancy on all of them is equivalent to constancy on every \(\approx_S\)-edge (an edge is witnessed by some shared \(y\); conversely each such set is pairwise \(\approx_S\)-connected through \(y\)), hence to \(\ker(q) \le \sigma^0(S)\) by connectivity, as in \hyperlink{stmt:22.1}{Theorem~22.1}. 2 and 4 are Fano and the DPI \autocite{cover2006}, the equality case in 4 being \hyperlink{stmt:22.6}{Theorem~22.6}(3). 3: chain rule plus nonnegativity of conditional mutual information; the Bayes-risk inequality holds because marginalizing away \(Y_w\) is a garbling, and garbling cannot lower Bayes risk (Blackwell--Sherman--Stein, \hyperlink{sec:20.1}{§20.1}).
\end{proof}

This is the promised enrichment of the Shannon base, and its shape deserves emphasis: each qualitative kernel principle acquires a quantitative counterpart ---

{\def\LTcaptype{none} % do not increment counter
\begin{longtable}[]{@{}
  >{\raggedright\arraybackslash}p{(\linewidth - 2\tabcolsep) * \real{0.4288}}
  >{\raggedright\arraybackslash}p{(\linewidth - 2\tabcolsep) * \real{0.5712}}@{}}
\toprule\noalign{}
\begin{minipage}[b]{\linewidth}\raggedright
kernel (exact)
\end{minipage} & \begin{minipage}[b]{\linewidth}\raggedright
graded / entropic
\end{minipage} \\
\midrule\noalign{}
\endhead
\bottomrule\noalign{}
\endlastfoot
\hyperlink{stmt:5.6}{Prop.~5.6}: answerable iff \(\ker(q) \le \sigma(S)\) & \hyperlink{stmt:22.1}{Thm~22.1} / 22.7(1): zero-error iff \(\ker(q)\le\sigma^0(S)\) iff \(H(q \mid Y_S) = 0\) \\
\hyperlink{stmt:5.5}{Def.~5.5}: unresolved pairs witness failure & \hyperlink{stmt:22.7}{Thm~22.7}(2): Fano converts residual entropy into forced error \\
\hyperlink{stmt:6.5}{Thm~6.5}: marginal value via separated pairs & \hyperlink{stmt:22.7}{Thm~22.7}(3): marginal value \(= I(Z; Y_w \mid Y_S)\); Bayes risk monotone \\
\hyperlink{stmt:14.2}{Lemma~14.2}: registration non-increasing & \hyperlink{stmt:22.7}{Thm~22.7}(4): DPI \\
\hyperlink{stmt:4.7}{Lemma~4.7}: canonical = most informative sound registration & \hyperlink{stmt:22.6}{Thm~22.6}(2--3): minimal sufficiency = coarsest lossless registration \\
\hyperlink{stmt:6.7}{Thm~6.7}: terminal consolidation \(D/\sigma(S)\) & \hyperlink{stmt:22.6}{Thm~22.6}: minimal sufficient statistic; deterministic case recovers 6.7 \\
\end{longtable}
}

--- so mutual information enters not as a replacement for the distinction-structure account but as its prior-weighted shadow, exactly as anticipated in Part I (§8, H3). The information bottleneck \autocite{tishby1999} now also finds its place: it is the \emph{lossy} relaxation of \hyperlink{stmt:22.6}{Theorem~22.6}---\hspace{0pt}minimize retained rate subject to retaining a stipulated amount of question-relevant information---\hspace{0pt}i.e., the graded form of budget-type (non-monotone) admissibility from \hyperlink{sec:16.2}{§16.2}.

\secn{23}{23}{Graded admissibility and the anatomy of synergy}

\subsecn{1}{23.1}{Graded ceilings}

\begin{thmdef}[\hypertarget{stmt:23.1}{Definition~23.1}]

A \textbf{graded admissibility structure} \(\widehat K\) on a coherent family assigns to each finite \(T\) a \textbf{ceiling experiment} \(\Gamma_{\widehat K}(T)\) with \(\Gamma_{\widehat K}(T) \preceq_B P_T\) (K1̂), with \(\Gamma_{\widehat K}(\varnothing)\) the null experiment (K0̂), the admissible readings on \(T\) being exactly the garblings of \(\Gamma_{\widehat K}(T)\) (K2̂: a principal Blackwell down-set). \(\widehat K\) is \textbf{monotone} if \(T \subseteq T' \Rightarrow \Gamma_{\widehat K}(T) \preceq_B \Gamma_{\widehat K}(T')\).
\end{thmdef}

As in \hyperlink{sec:16.2}{§16.2}, (K2̂) is a genuine restriction: natural resource constraints (rate limits, output budgets) carve out Blackwell down-sets that are not principal, and here the failure is aggravated by the absence of least upper bounds in the Blackwell order \autocite{bertschinger2014}---\hspace{0pt}there is in general no canonical completion of a non-principal admissible class to a ceiling. Part II's repair (a definite choice of ceiling within the class) remains available and remains a modeling act.

Part II's bracketing does \textbf{not} survive in graded form unamended, and establishing what does survive requires first fixing the object being compared.

\begin{thmdef}[\hypertarget{stmt:23.2}{Definition~23.2}\ {\normalfont\upshape(separable admissible reading)}]

For finite \(S\), a \textbf{separable admissible reading} of \(S\) under \(\widehat K\) is an experiment of the form
\[
\Big(\bigotimes_{v \in S} G_v\Big) \circ P_S,
\]
where each \(G_v\) is a channel on \(Y_v\), applied to the \(v\)-coordinate of the joint data, such that \(G_v \circ v\) is admissible on \(\{v\}\), i.e.~\(G_v \circ v \preceq_B \Gamma_{\widehat K}(\{v\})\). Each view is processed by its own admissible per-view decoder; no cross-view processing occurs.
\end{thmdef}

\begin{thmplain}[\hypertarget{stmt:23.3}{Proposition~23.3}\ {\normalfont\upshape(accumulation: monotonicity does not bound separable
readings)}]

There is a coherent CI family and a \emph{monotone} graded admissibility structure satisfying (K0̂)--(K2̂) with a separable admissible reading that strictly Blackwell-dominates the joint ceiling. Hence the bracketing of \hyperlink{stmt:14.6}{Proposition~14.6} fails in the graded theory without further axioms.
\end{thmplain}

\begin{proof}

Let \(D = \{0,1\}\) and let \(v, w\) be conditionally independent copies of the binary symmetric channel \(N\) with crossover \(\varepsilon \in (0, \tfrac12)\); the joint channel is \(P_{\{v,w\}} = N \otimes N\). Set \(\Gamma(\{v\}) = \Gamma(\{w\}) = \Gamma(\{v,w\}) = N\). This structure is monotone (all ceilings equal) and satisfies (K0̂)--(K2̂); in particular (K1̂) holds since \(N \preceq_B N \otimes N\) by marginalization. The identity processings \(G_v = G_w = \mathrm{id}\) are per-view admissible (\(\mathrm{id} \circ v = N \preceq_B N\)), and the resulting separable reading is \(N \otimes N\) itself. But \(N \otimes N \succ_B N\) strictly: dominance is marginalization; non-equivalence holds because the achievable posteriors differ---\hspace{0pt}under a uniform prior, one look yields posteriors in \(\{\varepsilon, 1-\varepsilon\}\), while two looks also realize the posterior \(\tfrac12\) (discordant outcomes) and \(\varepsilon^2/(\varepsilon^2 + (1-\varepsilon)^2)\) (concordant ones)---\hspace{0pt}and the concordant posterior is strictly below \(\varepsilon\). Every posterior from a garbling of one look is a convex combination of its posteriors and lies in \([\varepsilon,1-\varepsilon]\). Thus no garbling of \(N\) reproduces \(N\otimes N\).
\end{proof}

The mechanism deserves a name: \textbf{accumulation}. Independent noisy looks at the same candidate compound their evidence, so even \emph{re-reading the same view} is informative. In the exact theory this is invisible because combination is idempotent---\hspace{0pt}\(\ker(v) \vee \ker(v) = \ker(v)\), and a second reading of a deterministic view adds nothing. The failure of the Part II argument is therefore not a missing hypothesis but a category difference between exact combination (idempotent joins in a lattice) and graded combination (non-idempotent products of channels under couplings, in an experiment order where joins need not exist in general). Wherever ceilings are meant to cap \emph{total} extractive power, accumulation can carry separable readings past a joint ceiling that monotonicity alone leaves unguarded.

The bracket can be restored---\hspace{0pt}at a price that must be stated honestly.

\begin{thmdef}[\hypertarget{stmt:23.4}{Definition~23.4}\ {\normalfont\upshape(accumulation closure)}]

\(\widehat K\) is \textbf{accumulation-closed}---\hspace{0pt}axiom \textbf{(K3̂)}---\hspace{0pt}if for every finite \(S\), every separable admissible reading of \(S\) is a garbling of \(\Gamma_{\widehat K}(S)\).
\end{thmdef}

\begin{thmplain}[\hypertarget{stmt:23.5}{Proposition~23.5}\ {\normalfont\upshape(graded bracketing under (K3̂))}]

If \(\widehat K\) satisfies (K0̂)--(K3̂), then for every finite \(S\), prior \(\mu\), and question \(q\),
\[
e_\mu\big(q \mid \text{any separable admissible reading of } S\big)
\;\ge\;
e_\mu\big(q \mid \Gamma_{\widehat K}(S)\big)
\;\ge\;
e_\mu\big(q \mid P_S\big),
\]
the graded form of \hyperlink{stmt:14.6}{Proposition~14.6}'s bracket.
\end{thmplain}

\begin{proof}

Each separable reading is a garbling of \(\Gamma_{\widehat K}(S)\) by (K3̂), and \(\Gamma_{\widehat K}(S)\) is a garbling of \(P_S\) by (K1̂); garbling cannot lower Bayes risk (Blackwell--Sherman--Stein, \hyperlink{sec:20.1}{§20.1}).
\end{proof}

\begin{thmrem}[\hypertarget{stmt:23.6}{Remark~23.6}\ {\normalfont\upshape(the expressive trade-off)}]

The constant noisy ceiling of \hyperlink{stmt:23.3}{Proposition~23.3} models a total-output cap, but independently admissible one-copy readings together exceed it. Thus that particular cap and those local reading classes cannot simultaneously satisfy (K3̂). The honest joint ceiling \(\Gamma(S)=P_S\) does satisfy (K3̂), as does any realizable joint ceiling dominating the actual separable readings. This is an explicit conflict between a resource model and an accumulation axiom, not a proof that every budget model lacks a least completion. Adequacy of a proposed transfer is studied in \hyperlink{stmt:61.2}{Proposition~61.2} and Appendix D.11; full coupling fibers are governed by D.5. Down-set completion by itself supplies neither a physical coupling nor a principal achievable completion.
\end{thmrem}

\subsecn{2}{23.2}{Three mechanisms under one name}

Consider again the parity configuration of \hyperlink{sec:16.1}{§16.1}: \(D = \{00,01,10,11\}\), coordinate views \(v, w\), question \(q = b_1 \oplus b_2\), now with the uniform prior. Direct computation gives
\[\begin{gathered}I\big(q(Z); Y_v\big) = 0, \\ I\big(q(Z); Y_w\big) = 0, \\ I\big(q(Z); Y_v, Y_w\big) = 1 \text{ bit}:\end{gathered}\]
each coordinate is \emph{independent} of parity, yet the pair determines it. In the partial information decomposition (PID) literature this is the canonical example of \textbf{synergy}---\hspace{0pt}information carried only jointly \autocite{williams2010}. The present framework shows that at least three distinct mechanisms produce this signature, and separates them:

\begin{enumerate}
\def\labelenumi{\arabic{enumi}.}
\item
  \textbf{Prior-marginalization synergy} (full admissibility, exact views, a prior): the computation above. Note the contrast with the exact theory's verdict: by \hyperlink{stmt:6.5}{Theorem~6.5}, \(v\) alone \emph{does} strictly help \(q\) (it separates unresolved pairs such as \((00, 10)\)), while its mutual information with \(q\) is exactly zero. The discrepancy is not a paradox but a difference of accounting: the kernel counts \emph{distinctions drawn}, MI counts \emph{question-marginal uncertainty removed on prior-average}, and a view can do much of the first while doing none of the second. PID's ``synergy'' begins where these two ledgers diverge.
\item
  \textbf{Admissibility-induced synergy} (Part II, \hyperlink{sec:16.1}{§16.1}): with the invariance ceiling, even the \emph{distinction ledger} becomes superadditive---\hspace{0pt}\(\sigma^{\mathrm{sep}}_K = \bot\) while \(\sigma^{\mathrm{jnt}}_K = \ker(q)\)---\hspace{0pt}with no prior and no noise involved. Deterministic channels being a special case of graded ones, this mechanism persists verbatim in the graded theory.
\item
  \textbf{Support-geometric synergy} (Proposition 22.3): zero-error content is superadditive purely through the intersection pattern of noise supports, with full admissibility and before any prior enters.
\end{enumerate}

The examples distinguish changes in prior-relative information, admissible resolving power, and support geometry. A general axiomatic decomposition connecting these mechanisms is not established by those examples. Part VII gives exact-chain evaluations and coupling-fiber excess as separate constructions, and Appendix D.4 proves how the latter changes under quotient averaging. Neither construction alone proves a PID impossibility theorem.

A closing note on \hyperlink{sec:23.1}{§23.1}'s accumulation phenomenon, which stands to these three mechanisms as their reverse. Where each of (1)--(3) makes the whole read \emph{more} than the sum of its parts, accumulation makes separable readings exceed a joint \emph{ceiling}---\hspace{0pt}an anti-synergetic divergence tied not to the joint distribution but to the multiplicity of admissible readings and the choice of ceiling. PID cannot see it even in principle: PID decomposes a \emph{fixed} joint distribution of sources and target, so mechanisms that vary with the ceiling assignment or with how many admissible looks a decoder takes lie outside its object of study altogether. Accumulation is related to, but distinct from, PID's redundancy: redundant information is carried by each source separately \emph{within} one fixed joint; accumulated information is created by compounding looks \emph{across} readings that no single fixed joint represents. It is recorded here as a fourth entry in the anatomy---\hspace{0pt}the one that only the admissibility-plus-grading combination exposes.

\secn{24}{24}{The zero-noise limit: recovery of the kernel}

\begin{thmplain}[\hypertarget{stmt:24.1}{Lemma~24.1}\ {\normalfont\upshape(deterministic Blackwell order is view refinement)}]

For deterministic channels \(v, v'\) on \(D\): \(v \succeq_B v'\) iff \(v' \preceq v\) in the refinement preorder of \hyperlink{stmt:4.3}{Definition~4.3}, iff \(\ker(v') \le \ker(v)\).
\end{thmplain}

\begin{proof}

(\(\Leftarrow\)) \hyperlink{stmt:4.4}{Lemma~4.4} provides a deterministic mediating map, which is in particular a channel. (\(\Rightarrow\)) Let \(v' = G \circ v\) with \(G\) a channel. For each \(Z\), \(\delta_{v'(Z)} = G(\cdot \mid v(Z))\), so \(G(\cdot \mid v(Z))\) is the point mass at \(v'(Z)\); if \(v(Z) = v(Z_1)\) the two point masses coincide, so \(v'(Z) = v'(Z_1)\), i.e.~\(\ker(v') \le \ker(v)\), and \hyperlink{stmt:4.4}{Lemma~4.4} converts this into a deterministic factorization.
\end{proof}

\begin{thmplain}[\hypertarget{stmt:24.2}{Theorem~24.2}\ {\normalfont\upshape(recovery)}]

Restrict a coherent family to deterministic channels (equivalently, let all noise vanish). Then:

\begin{enumerate}
\def\labelenumi{\arabic{enumi}.}
\tightlist
\item
  couplings are unique, so the family is determined by its single views, and \hyperlink{stmt:21.3}{Remark~21.3}'s indeterminacy disappears;
\item
  the Blackwell preorder restricts to the refinement preorder, on which suprema exist and are computed in \(\mathrm{Part}(D)\) (Lemma 24.1, \hyperlink{stmt:3.3}{Fact~3.3});
\item
  \(\sigma^{0}(S) = \sigma^{=}(S) = \sigma(S)\) (Lemma 21.7), and \hyperlink{stmt:22.1}{Theorems~22.1} and 22.2 both reduce to \hyperlink{stmt:5.6}{Proposition~5.6};
\item
  the inequalities of \hyperlink{stmt:22.3}{Propositions~22.3} and 22.4 are equalities (Theorem 6.2(1))---\hspace{0pt}deterministic couplings are unique, so every family is trivially CI and both superadditivity mechanisms vanish; the adjunction T1 holds;
\item
  the minimal sufficient statistic is the terminal consolidation \(D/\sigma(S)\) (Theorem 22.6(4) = \hyperlink{stmt:6.7}{Theorem~6.7});
\item
  the entropic criteria specialize: \(H(q(Z) \mid R_S(Z)) = 0\) for full-support priors iff \(\ker(q) \le \sigma(S)\), and the DPI holds with equality precisely for registrations preserving \(\ker\)-structure, recovering \hyperlink{stmt:14.2}{Lemma~14.2}'s equality case.
\end{enumerate}
\end{thmplain}

\begin{proof}

Each item collects the deterministic specializations established in the lemma or theorem cited within it; no new argument is required.
\end{proof}

The exact theory of Parts I--II is thus not an idealization discarded by the graded theory but its precise zero-noise fiber---\hspace{0pt}the sublocus on which couplings are canonical, the content order is a complete lattice, the three strata coincide, and every graded theorem collapses onto its kernel ancestor.

\subsubsection{The structural collapse}\label{the-structural-collapse}

The recovery theorem invites a reading of Parts I--III as a single object: a record of how much algebraic structure the notion of content retains as the standing assumptions are relaxed one by one. Two things degrade, and they degrade independently, so the record must track both: the \textbf{codomain} in which content lives (is it still a lattice?), and the \textbf{compositionality of the content map} over unions of view sets (does content over \(S \cup S'\) still decompose as a join of the contents over the parts?). The table collects the verdicts established across the three parts.

{\def\LTcaptype{none} % do not increment counter
\begin{longtable}[]{@{}
  >{\raggedright\arraybackslash}p{(\linewidth - 6\tabcolsep) * \real{0.1522}}
  >{\raggedright\arraybackslash}p{(\linewidth - 6\tabcolsep) * \real{0.1488}}
  >{\raggedright\arraybackslash}p{(\linewidth - 6\tabcolsep) * \real{0.2324}}
  >{\raggedright\arraybackslash}p{(\linewidth - 6\tabcolsep) * \real{0.4666}}@{}}
\toprule\noalign{}
\begin{minipage}[b]{\linewidth}\raggedright
layer
\end{minipage} & \begin{minipage}[b]{\linewidth}\raggedright
content object
\end{minipage} & \begin{minipage}[b]{\linewidth}\raggedright
codomain
\end{minipage} & \begin{minipage}[b]{\linewidth}\raggedright
content map over unions
\end{minipage} \\
\midrule\noalign{}
\endhead
\bottomrule\noalign{}
\endlastfoot
kernel (Part I) & \(\sigma(S)\) & complete lattice \(\mathrm{Part}(D)\) & preserves joins; full adjunction (Thm 6.2) \\
separable admissibility (Part II) & \(\sigma^{\mathrm{sep}}_K(S)\) & complete lattice \(\mathrm{Part}(D)\) & preserves joins; full adjunction (Thm 15.1) \\
joint admissibility (Part II) & \(\sigma^{\mathrm{jnt}}_K(S)\) & complete lattice \(\mathrm{Part}(D)\) & superadditive at non-separable \(K\) (Prop. 15.2); inverted at non-monotone \(K\) (§16.2); adjunction iff separable (Cor. 15.4) \\
graded, zero-error stratum & \(\sigma^{0}(S)\) (partition shadow of \(\sigma^{G}\)) & complete lattice \(\mathrm{Part}(D)\) & superadditive even on CI families (Prop. 22.3); adjunction restored at the graph level, where CI content composes exactly (Prop. 22.3′)---\hspace{0pt}the loss is the reflection to partitions \\
graded, statistical stratum & \(\sigma^{=}(S)\) & complete lattice \(\mathrm{Part}(D)\) & superadditive exactly through coupling; joins and adjunction restored on CI families (Prop. 22.4) \\
graded, full stratum & \(\widehat{\sigma}(S)\) & Blackwell preorder; binary-state lattice, joins may fail in general \autocite{bertschinger2014} & ill-posed: the union's content depends on the coupling, which the parts do not determine (Prop. 22.5, Rem. 21.3) \\
graded admissibility (§23) & readings under \(\widehat K\) & Blackwell preorder & separable/joint bracket only under accumulation closure (K3̂), at the cost of budget semantics (Rem. 23.6) \\
\end{longtable}
}

The partition-valued strata keep a lattice codomain even when their content maps fail to preserve joins. The full experiment order adds a different issue: in general some joins fail to exist, and even where they exist a supplied joint experiment can be strictly above the join. These distinctions motivate separate exact and decision-theoretic comparisons; they do not make every noisy experiment a nonlattice case.

\subsubsection{The content-system schema}\label{the-content-system-schema}

The collapse table invites, and the accumulated results now permit, a single definition of which every content notion in Parts I--III is an instance. Let \(\mathcal{F}(V)\) denote the finite subsets of \(V\), a join-semilattice under union.

\begin{thmdef}[\hypertarget{stmt:24.3}{Definition~24.3}\ {\normalfont\upshape(content system)}]

A \textbf{content system} over \((D, V)\) is a monotone map \(\mathcal{C} : \mathcal{F}(V) \to \mathsf{C}\) into a preordered \textbf{content codomain} \(\mathsf{C}\), together with the \textbf{comparison cells}
\[
\mathcal{C}(T) \vee \mathcal{C}(T') \;\le\; \mathcal{C}(T \cup T'),
\]
wherever the join on the left exists---\hspace{0pt}i.e.~a lax morphism of join-semilattices. \(\mathcal{C}\) is \textbf{strict} if all cells exist and are equalities. The non-invertible cells are the system's \textbf{anomalies}.
\end{thmdef}

The manuscript's content maps are then classified by two coordinates---\hspace{0pt}the codomain and the status of the cells:

{\def\LTcaptype{none} % do not increment counter
\begin{longtable}[]{@{}
  >{\raggedright\arraybackslash}p{(\linewidth - 4\tabcolsep) * \real{0.2315}}
  >{\raggedright\arraybackslash}p{(\linewidth - 4\tabcolsep) * \real{0.2373}}
  >{\raggedright\arraybackslash}p{(\linewidth - 4\tabcolsep) * \real{0.5311}}@{}}
\toprule\noalign{}
\begin{minipage}[b]{\linewidth}\raggedright
content map
\end{minipage} & \begin{minipage}[b]{\linewidth}\raggedright
codomain \(\mathsf{C}\)
\end{minipage} & \begin{minipage}[b]{\linewidth}\raggedright
comparison cells
\end{minipage} \\
\midrule\noalign{}
\endhead
\bottomrule\noalign{}
\endlastfoot
\(\sigma\) (Part I) & \(\mathrm{Part}(D)\) & strict (Theorem 6.2) \\
\(\sigma^{\mathrm{sep}}_K\) & \(\mathrm{Part}(D)\) & strict (Theorem 15.1) \\
\(\sigma^{\mathrm{jnt}}_K\), monotone \(K\) & \(\mathrm{Part}(D)\) & strict iff \(K\) separable (Corollary 15.4); \(\Delta_K\) is the non-invertible cell \\
\(\sigma^{G}\) on CI families & tolerance lattice \(\mathrm{Tol}(D)^{\sqsubseteq}\) & strict (Proposition 22.3′) \\
\(\sigma^{0}\) & \(\mathrm{Part}(D)\) & lax even under CI, through reflection (Proposition 22.3) \\
\(\sigma^{=}\) & \(\mathrm{Part}(D)\) & strict on CI families; lax through coupling (Proposition 22.4) \\
\(\widehat{\sigma}\) & Blackwell preorder & no general join-based cell; supplied coupling remains additional data (Proposition 22.5) \\
\end{longtable}
}

Three facts hold at the level of the schema itself.

\begin{thmplain}[\hypertarget{stmt:24.4}{Proposition~24.4}\ {\normalfont\upshape(adjunction schema)}]

Let \(\mathsf C\) be a complete lattice and let \(\mathcal C:\mathcal P(V)\to\mathsf C\) satisfy \(\mathcal C(T)=\bigvee_{v\in T}\mathcal C(\{v\})\), including \(\mathcal C(\varnothing)=\bot\). Then \(\mathcal C\) is left adjoint to \(\tau(\pi)=\{v:\mathcal C(\{v\})\le\pi\}\). In particular this applies to a grounded singleton-generated system on finite subsets when \(V\) is finite, or to its arbitrary-join extension when \(V\) is infinite.
\end{thmplain}

\begin{proof}

\(\mathcal C(T)\le\pi\) iff each singleton generator is below \(\pi\), iff \(T\subseteq\tau(\pi)\). The closure and fixed-point consequences are \hyperlink{stmt:3.4}{Fact~3.4}.
\end{proof}

\begin{thmdef}[\hypertarget{stmt:24.5}{Definition~24.5}\ {\normalfont\upshape(gluing)}]

Equip \(\mathcal{F}(V)\) with the \textbf{positive union coverage}: a nonempty family \(\{T_i\}_i\) covers \(T\) iff \(\bigcup_i T_i = T\). Grounding at \(\varnothing\) is imposed separately. This convention is needed for the witness presheaf of \hyperlink{sec:38}{§38}. A content system \textbf{glues} if its comparison cells are invertible on covers: \(\bigvee_i \mathcal{C}(T_i) = \mathcal{C}(T)\) whenever \(\{T_i\}\) covers \(T\). Separability (Definition 15.3) is exactly gluing for \(\sigma^{\mathrm{jnt}}_K\); \hyperlink{stmt:6.3}{Propositions~6.3}, 22.3′, and 22.4(3) are gluing statements; and the sheaf-theoretic extension H4 is, in this vocabulary, the obstruction theory of non-invertible cells, with the interval-typed anomaly (Definition 15.3) supplying composable coefficients. This covariant join equation is the definition used here; it is distinct from the contravariant Set-valued sheaf axiom. That obstruction theory is executed in Part V (Theorems 37.3, 37.5, 38.4, 39.2, 41.2, 42.1), with \textbf{descent} installed as a tracked coordinate of the architecture alongside the collapse table's two, representability (Remark 24.6), and fidelity (Part IV, \hyperlink{sec:32.2}{§32.2}).
\end{thmdef}

\begin{thmrem}[\hypertarget{stmt:24.6}{Remark~24.6}\ {\normalfont\upshape(down-set completion and its limits)}]

For an information preorder \(P\), let \(\operatorname{Down}(P)\) be all lower sets, ordered by inclusion. Intersections and unions give a complete lattice and \(p\mapsto\downarrow p\) is an order embedding modulo equivalence. These lower sets need not be directed ideals. A lower set represents available alternative experiments, not their assembled joint experiment. In particular \(\downarrow a\cup\downarrow b\) can be strictly smaller than \(\downarrow(a\vee b)\) even when the latter join exists.

For a \textbf{specified} closure rule on lower-set assignments that is preserved by intersections, intersections of closed majorants yield a least closed assignment. This formal fact does not prove that graded accumulation closure for arbitrary supplied couplings has such a rule, that the assignment is principal, or that its generator is achievable. Those are additional obligations. The Blackwell no-join obstruction is general, with the binary-state exception of \hyperlink{stmt:22.5}{Proposition~22.5}. Appendix D.5 treats unconstrained coupling minima directly.
\end{thmrem}

\begin{thmrem}[\hypertarget{stmt:24.7}{Remark~24.7}\ {\normalfont\upshape(quantitative cells)}]

Le Cam deficiency (§25) makes experiments a generalized metric space in the sense of enriched category theory \autocite{lawvere1973,lecam1964}, and an ideal-valued content system a metrically enriched one. The approximate-gluing programme flagged for H4 is then the study of comparison cells invertible up to a stated deficiency---\hspace{0pt}a quantitative laxity---\hspace{0pt}rather than a separate theory.
\end{thmrem}

\secn{25}{25}{Established results and scope}

The finite answerability and sufficiency results establish the three information strata. Accumulation explains why monotone local ceilings alone need not bound joint readings; Appendix D.3 characterizes precisely when independent replication is idempotent. Down-set completion describes alternatives but does not preserve all existing joins or choose a physical coupling. Appendix A gives one finite-state Borel extension and states the additional hypotheses needed elsewhere.
